% ============================================================
\chapter{Introduction to Ricci Flow}
\label{ch:ricci_flow}
% ============================================================

\section{Motivation and definition}

The Ricci flow, introduced by R.~Hamilton in 1982, is a fundamental tool in
differential geometry. The idea is to deform a Riemannian metric in the direction
of its Ricci curvature, thereby smoothing the geometry over time.

\begin{intuition}
The Ricci flow is the geometric analogue of the heat equation: just as heat
diffuses to equalize temperature, the Ricci flow ``diffuses'' curvature to make
it more uniform. Regions of large positive curvature contract, while regions of
negative curvature expand.
\end{intuition}

\begin{definition}[Ricci flow]
\index{Ricci flow}
Let $(M, g_0)$ be a compact Riemannian manifold. The \emph{Ricci flow} is the
evolution equation:
\[
    \frac{\partial g}{\partial t} = -2\operatorname{Ric}(g), \qquad g(0) = g_0.
\]
\end{definition}

\begin{remark}
The factor $-2$ is a convention. The negative sign ensures that regions of positive
Ricci curvature contract (as physically expected). The equation is a quasilinear
parabolic system for the metric $g_{ij}$.
\end{remark}

\begin{definition}[Normalized Ricci flow]
\index{normalized Ricci flow}
The \emph{normalized Ricci flow} preserves the volume:
\[
    \frac{\partial g}{\partial t} = -2\operatorname{Ric}(g) + \frac{2}{n}\bar{R}\, g,
\]
where $\bar{R} = \frac{\int_M R\, dV_g}{\int_M dV_g}$ is the average scalar curvature
and $n = \dim M$.
\end{definition}

\section{Short-time existence and uniqueness}

\begin{theorem}[Hamilton, 1982]
\label{thm:hamilton_existence}
\index{Hamilton's existence theorem}
Let $(M, g_0)$ be a compact Riemannian manifold. There exists $T > 0$ and a unique
smooth family of metrics $(g(t))_{t \in [0, T)}$ satisfying the Ricci flow with
initial condition $g(0) = g_0$.
\end{theorem}

\begin{remark}
Hamilton's original proof uses the Nash-Moser theorem. DeTurck (1983) subsequently
gave a simplified proof by introducing a modification of the flow that makes it
strictly parabolic.
\end{remark}

\begin{definition}[DeTurck trick]
\index{DeTurck trick}
Let $\tilde{g}$ be a fixed reference metric on $M$. The \emph{Ricci-DeTurck flow} is:
\[
    \frac{\partial g}{\partial t} = -2\operatorname{Ric}(g)
    + \mathcal{L}_{W(g, \tilde{g})} g,
\]
where $W$ is a vector field depending on $g$ and $\tilde{g}$ (defined via the
difference of Christoffel symbols). This flow is strictly parabolic.
\end{definition}

\begin{proposition}
The Ricci-DeTurck flow is equivalent to the Ricci flow modulo diffeomorphisms:
if $g(t)$ solves the Ricci-DeTurck flow, there exists a family of diffeomorphisms
$\varphi_t : M \to M$ such that $\varphi_t^* g(t)$ solves the Ricci flow.
\end{proposition}

\section{Evolution of curvature}

\begin{theorem}[Evolution of scalar curvature]
\label{thm:evolution_R}
Under the Ricci flow, the scalar curvature evolves by:
\[
    \frac{\partial R}{\partial t} = \Delta R + 2\abs{\operatorname{Ric}}^2.
\]
\end{theorem}

\begin{proof}
In local coordinates, $R = g^{ij} R_{ij}$. Differentiating:
\[
    \frac{\partial R}{\partial t} = g^{ij} \frac{\partial R_{ij}}{\partial t}
    + \frac{\partial g^{ij}}{\partial t} R_{ij}.
\]
Since $\frac{\partial g^{ij}}{\partial t} = 2R^{ij}$, and using the evolution
formula for the Ricci tensor:
\[
    \frac{\partial R_{ij}}{\partial t} = \Delta R_{ij} + 2R_{ikjl}R^{kl}
    - 2R_{ik}R_j^{\;k},
\]
contraction yields the stated result.
\end{proof}

\begin{corollary}[Maximum principle for $R$]
If the initial scalar curvature satisfies $R(g_0) \geq R_{\min}$, then for all
$t \in [0, T)$:
\[
    R(g(t)) \geq \frac{R_{\min}}{1 - \frac{2R_{\min}}{n} t}.
\]
In particular, if $R_{\min} > 0$, the flow develops a singularity in finite time
$T \leq \frac{n}{2R_{\min}}$.
\end{corollary}

\begin{theorem}[Evolution of the Ricci tensor]
Under the Ricci flow:
\[
    \frac{\partial R_{ij}}{\partial t} = \Delta R_{ij}
    + 2R_{ikjl}R^{kl} - 2R_{ik}R_j^{\;k}.
\]
\end{theorem}

\begin{theorem}[Evolution of the Riemann tensor]
Under the Ricci flow:
\[
    \frac{\partial R_{ijkl}}{\partial t} = \Delta R_{ijkl}
    + 2(B_{ijkl} - B_{ijlk} + B_{ikjl} - B_{iljk}),
\]
where $B_{ijkl} = R_{ipjq}R^{p\phantom{k}q}_{\phantom{p}k\phantom{q}l}$ (summing
over repeated indices with the Einstein convention).
\end{theorem}

\section{Hamilton's maximum principle}

\begin{theorem}[Scalar maximum principle]
\label{thm:scalar_maximum}
\index{maximum principle}
Let $u : M \times [0, T) \to \R$ satisfy:
\[
    \frac{\partial u}{\partial t} \geq \Delta_{g(t)} u + F(u),
\]
where $g(t)$ solves the Ricci flow and $F$ is locally Lipschitz. If
$u(x, 0) \geq c$ for all $x \in M$, then $u(x, t) \geq \phi(t)$ where $\phi$
solves the ODE $\phi' = F(\phi)$, $\phi(0) = c$.
\end{theorem}

\begin{proof}[Proof sketch]
One uses the parabolic comparison principle. Let $x_0(t)$ be a point where
$u(\cdot, t)$ attains its spatial minimum. At this point $\Delta_{g(t)} u \geq 0$
(since it is a minimum), so:
\[
    \frac{\partial u}{\partial t}\bigg|_{x_0(t)} \geq F(u(x_0(t), t)).
\]
It follows (via a regularization argument, since $x_0(t)$ may jump) that
$u_{\min}(t) = \min_M u(\cdot, t)$ satisfies $\frac{d}{dt} u_{\min} \geq F(u_{\min})$
in the barrier sense. Comparison with the ODE $\phi' = F(\phi)$, $\phi(0) = c$
yields $u_{\min}(t) \geq \phi(t)$. The typical application is $u = R$ (scalar
curvature), $F(u) = \frac{2}{n}u^2$, giving the bound
$R \geq \frac{R_{\min}}{1 - \frac{2R_{\min}}{n}t}$.
\end{proof}

\begin{warning}[title=\textbf{The metric depends on time}]
The Laplacian $\Delta_{g(t)}$ depends on the evolving metric. The maximum principle
therefore requires careful treatment, since metric balls and norms change with time.
\end{warning}

\begin{theorem}[Hamilton's tensor maximum principle]
\label{thm:tensor_maximum}
Let $T_{ij}$ be a symmetric tensor evolving by:
\[
    \frac{\partial T_{ij}}{\partial t} = \Delta T_{ij} + N_{ij}(T, g),
\]
where $N_{ij}$ satisfies: if $T_{ij} \geq 0$ and $T_{ij}v^j = 0$ at a point,
then $N_{ij} v^i v^j \geq 0$. If $T_{ij}(0) \geq 0$, then $T_{ij}(t) \geq 0$
for all $t \in [0, T)$.
\end{theorem}

\begin{proof}[Proof sketch]
Consider the smallest eigenvalue $\lambda_{\min}(t)$ of $T_{ij}$ over $M$. If
$\lambda_{\min}(t_0) = 0$ at a point $x_0$ with eigenvector $v$, then
$T_{ij}v^j = 0$ at that point, so $N_{ij}v^i v^j \geq 0$ by hypothesis. A maximum
principle argument (applied to the quantity $T_{ij}v^i v^j$ along curves of parallel
vectors) shows that $\lambda_{\min}$ cannot cross zero from above. The technical
difficulty is that the eigenvector depends on the point; this is handled by working
on the orthonormal frame bundle.
See \textsc{Hamilton}, \emph{Three-manifolds with positive Ricci curvature},
J.~Differential Geom.~17 (1982), Theorem~9.1.
\end{proof}

\begin{corollary}[Preservation of nonnegative Ricci]
If $\operatorname{Ric}(g_0) \geq 0$, then $\operatorname{Ric}(g(t)) \geq 0$
for all $t$ in the interval of existence.
\end{corollary}

\section{Fundamental examples}

\begin{example}[Sphere $S^n$]
\index{Ricci flow!sphere}
On the sphere $(S^n, g_0)$ with the round metric of sectional curvature $1$, we have
$\operatorname{Ric}(g_0) = (n-1)g_0$. The Ricci flow has the solution:
\[
    g(t) = (1 - 2(n-1)t)\, g_0.
\]
The metric shrinks uniformly and vanishes at $T = \frac{1}{2(n-1)}$. The sphere
contracts to a point in finite time.
\end{example}

\begin{example}[Hyperbolic space]
\index{Ricci flow!hyperbolic space}
On $(\mathbb{H}^n, g_0)$ with $\operatorname{Ric}(g_0) = -(n-1)g_0$:
\[
    g(t) = (1 + 2(n-1)t)\, g_0.
\]
The metric expands uniformly. The flow exists for all $t > 0$.
\end{example}

\begin{example}[Ricci solitons]
\index{Ricci soliton}
A \emph{Ricci soliton} is a solution that evolves only by diffeomorphisms and
rescaling:
\[
    \operatorname{Ric}(g) + \mathcal{L}_X g + \lambda g = 0,
\]
where $X$ is a vector field and $\lambda \in \R$. The soliton is called
shrinking ($\lambda > 0$), steady ($\lambda = 0$), or expanding ($\lambda < 0$).
\end{example}

\begin{example}[Flow on surfaces]
\index{Ricci flow!surfaces}
In dimension $2$, $\operatorname{Ric} = \frac{R}{2}g$ and the Ricci flow becomes:
\[
    \frac{\partial g}{\partial t} = -Rg.
\]
The normalized Ricci flow is $\frac{\partial g}{\partial t} = (\bar{R} - R)g$.
Hamilton (1988) and Chow (1991) showed that on any compact surface, the normalized
flow converges to a metric of constant curvature.
\end{example}

\section{Convergence in dimension 3}

\begin{theorem}[Hamilton, 1982]
\label{thm:hamilton_3d}
Let $(M^3, g_0)$ be a compact $3$-dimensional Riemannian manifold with strictly
positive Ricci curvature. The normalized Ricci flow converges to a metric of
constant positive sectional curvature. In particular, $M$ is diffeomorphic to a
quotient $S^3/\Gamma$.
\end{theorem}

\begin{remark}
This foundational result motivated the Hamilton-Perelman program, culminating in
Perelman's proof of the Poincar\'e conjecture (2002--2003) in dimension $3$.
\end{remark}

\begin{keyformulas}[title=\textbf{Evolution equations under the Ricci flow}]
\begin{align*}
    \frac{\partial g_{ij}}{\partial t} &= -2R_{ij} \\[6pt]
    \frac{\partial}{\partial t}\Gamma_{ij}^k &= -g^{kl}(\nabla_i R_{jl}
        + \nabla_j R_{il} - \nabla_l R_{ij}) \\[6pt]
    \frac{\partial R}{\partial t} &= \Delta R + 2\abs{\operatorname{Ric}}^2 \\[6pt]
    \frac{\partial R_{ij}}{\partial t} &= \Delta R_{ij} + 2R_{ikjl}R^{kl}
        - 2R_{ik}R_j^{\;k} \\[6pt]
    \frac{\partial}{\partial t}dV_g &= -R\, dV_g
\end{align*}
\end{keyformulas}

\section{Exercises}

\begin{exercise}
Verify that $g(t) = (1 - 2(n-1)t) g_0$ solves the Ricci flow on $(S^n, g_0)$.
Compute the extinction time and $\operatorname{diam}(S^n, g(t))$ as a function of $t$.
\end{exercise}

\begin{exercise}
Show that under the Ricci flow, the volume evolves by:
$\frac{d}{dt}\operatorname{Vol}(M, g(t)) = -\int_M R\, dV_g$.
Deduce that the volume decreases when $R > 0$.
\end{exercise}

\begin{exercise}
Let $(M^2, g_0)$ be a compact surface. Show that the normalized Ricci flow preserves
total area and that the Euler characteristic $\chi(M)$ determines the sign of
$\bar{R}$ via the Gauss-Bonnet theorem.
\end{exercise}

\begin{exercise}[Hamilton's cigar soliton]
Show that the metric on $\R^2$:
\[
    g = \frac{dx^2 + dy^2}{1 + x^2 + y^2}
\]
is a steady Ricci soliton (the ``cigar soliton''). Compute its Gaussian curvature.
\end{exercise}

\begin{exercise}
Let $g(t) = \varphi(t) g_0$ be a solution of the Ricci flow with $\varphi(0) = 1$.
Show that this is only possible if $(M, g_0)$ is an Einstein manifold:
$\operatorname{Ric}(g_0) = \lambda g_0$. Determine $\varphi(t)$ as a function of
$\lambda$.
\end{exercise}

\begin{exercise}
Using the maximum principle, show that if $R(g_0) \geq R_{\min} > 0$ on a compact
manifold, then $R(g(t)) \to +\infty$ as $t$ approaches the maximal existence time.
\end{exercise}
